\documentclass[border=10pt]{standalone}
\usepackage{fontspec}
\setmainfont{Noto Serif}[Ligatures=TeXOff]
\XeTeXgenerateactualtext=1
\newfontfamily\bengalifont{Noto Serif Bengali}[Script=Bengali, NFSSFamily=bnfam]
\usepackage[Latin,Bengali]{ucharclasses}
% group-free font switching (as in header-common.tex): robust inside braces, cells and quotes
\makeatletter
\def\BN@saved{\rmdefault}
\def\BN@bn{bnfam}
\newcommand\BNin{\edef\BN@tmp{\f@family}\ifx\BN@tmp\BN@bn\else\let\BN@saved\BN@tmp\fi\fontfamily{bnfam}\selectfont}
\newcommand\BNout{\fontfamily{\BN@saved}\selectfont}
\makeatother
\setTransitionsFor{Bengali}{\BNin}{\BNout}
% Typographic quotes and dashes are not in a ucharclasses class (class 0), so a Bangla word right
% after “ or ‘ would not trigger the Bangla font. Treat them as Latin so that every switch fires.
\XeTeXcharclass"2018=\XeTeXcharclass`A
\XeTeXcharclass"2019=\XeTeXcharclass`A
\XeTeXcharclass"201C=\XeTeXcharclass`A
\XeTeXcharclass"201D=\XeTeXcharclass`A
\XeTeXcharclass"2013=\XeTeXcharclass`A
\XeTeXcharclass"2014=\XeTeXcharclass`A
\XeTeXcharclass"2026=\XeTeXcharclass`A
\usepackage{newunicodechar}
\newunicodechar{।}{{\bengalifont\char"0964}}
\newunicodechar{॥}{{\bengalifont\char"0965}}

\usepackage[table]{xcolor}
\usepackage{booktabs,tabularx,array}
\definecolor{bapilrule}{HTML}{1F2937}
\definecolor{bapilstripe}{HTML}{F3F4F6}
\renewcommand{\arraystretch}{1.28}
\setlength{\tabcolsep}{5pt}

\begin{document}
\begin{minipage}{17cm}
{\large\bfseries সারণি ৪৭।}\ {\large কৌশলগত দৃশ্যকল্প ও করণীয়}\par\vspace{7pt}
\small
\rowcolors{2}{bapilstripe}{white}
\begin{tabularx}{\linewidth}{@{}>{\hsize=0.681\hsize\raggedright\arraybackslash}X>{\hsize=1.106\hsize\raggedright\arraybackslash}X>{\hsize=0.979\hsize\raggedright\arraybackslash}X>{\hsize=1.234\hsize\raggedright\arraybackslash}X@{}}
\arrayrulecolor{bapilrule}\toprule
\rowcolor{white}\textbf{দৃশ্যকল্প} & \textbf{বিবরণ} & \textbf{প্রাথমিক সংকেত} & \textbf{কৌশলগত করণীয়} \\
\midrule
\textbf{ত্বরান্বিত} & শক্তিশালী রাজনৈতিক পৃষ্ঠপোষকতা, জাতীয় কম্পিউটিং ও ডেটা-বিনিময় আগেভাগে উপলব্ধ & অর্থায়নসহ এআই নীতি গৃহীত; NRDEX চালু; সংস্থাগুলোর সরঞ্জাম চাহিদা & H2-এর গ্রহণ এগিয়ে আনা; আরও মন্ত্রণালয়ে সম্প্রসারণ; গতির কারণে সুরক্ষা যেন না কমে তাই গেট অপরিবর্তিত \\
\textbf{ভিত্তি} & ডেটা প্রবেশাধিকার ও অর্থায়নে স্বাভাবিক বিলম্বসহ স্থির অগ্রগতি & প্রায় দুই ত্রৈমাসিকের মধ্যে মাইলফলক অর্জিত & রোডম্যাপ অনুসরণ; পর্যালোচনার মাধ্যমে পরিধি সমন্বয় \\
\textbf{সীমিত} & অর্থায়ন হ্রাস, ধীর ডেটা-চুক্তি বা পৃষ্ঠপোষকতা হারানো & MS2 বা MS4 অনর্জিত; বাজেট বিলম্ব; কর্মী বদল & মূল অংশ রক্ষা (P1, বেঞ্চমার্ক, শাসন); P3, P5 ও P6 স্থগিত; গবেষণা অনুদান সন্ধান; যা আছে তা প্রকাশ \\
\textbf{প্রযুক্তিগত পরিবর্তন} & অনেক শক্তিশালী বা সস্তা মডেল সম্ভাবনার পরিধি বদলে দেয় & অভিযোজন ছাড়াই উন্মুক্ত মডেল সীমা অতিক্রম করে; ব্যয় দ্রুত কমে & মডেল নির্বাচন পুনরায় (সারণি ২৯); অপ্রয়োজনীয় BPFM পর্যায় বাদ; মূল্যায়ন ও ডেটায় বিনিয়োগ \\
\textbf{শাসনগত ধাক্কা} & বাংলাদেশে বা বিদেশে গুরুতর এআই ঘটনায় জন-আস্থা হ্রাস & সংবাদমাধ্যমে প্রচার; নতুন বিধিনিষেধ; অভিযোগ & প্রভাবিত ব্যবস্থা স্থগিত; স্বাধীন পর্যালোচনা; ফলাফল প্রকাশ; নিয়ন্ত্রণ জোরদার; উন্মুক্ত যোগাযোগ \\
\bottomrule
\end{tabularx}
\par\vspace{3pt}{\scriptsize\color{gray} উৎস: BAPIL টেকনিক্যাল ব্লুপ্রিন্ট, সংস্করণ ১.০ (২০২৬)। CC BY 4.0।}
\end{minipage}
\end{document}
